\documentclass[11pt]{article}
\usepackage[margin=1in]{geometry}
\usepackage{hyperref}
\hypersetup{colorlinks=true,linkcolor=blue,urlcolor=blue}
\pagestyle{empty}

\begin{document}

\begin{center}
{\large HyperTensor Papers I--X}

\vspace{8pt}

William Ken Ohara Stewart

NagusameCS Independent Research

\texttt{nagusamecs@gmail.com}

\url{https://github.com/NagusameCS/HyperTensor}

\vspace{8pt}

May 2026
\end{center}

\vspace{12pt}

Dear Reader,

Thank you for taking the time to look through this work. My name is William
Ken Ohara Stewart. I am eighteen years old and a high school student in the
United States. I am not a professional researcher, I do not have a lab or an
advisor, and I am certain there are mistakes in these pages that a trained eye
would catch immediately. I ask for your patience with those. Everything here
was done on a laptop with an RTX 4070, a rented cloud GPU when a bigger one was
needed, and a lot of late nights.

What follows is a set of ten papers that try to answer one question: can you
compress a language model using only the geometry of its own weights, with no
training data and no calibration, and still have it work? The answer, as of
Paper~X, is a qualified yes---but the qualification matters more than the yes
itself. Below is a short summary of each part so you know what you are getting
into.

\vspace{8pt}

\noindentPaper~I --- Geodesic Runtime Compression.}
The core measurement. We compress only the attention weights (Q, K, V) of
Llama-3.1-8B by projecting them onto their own principal components, with no
calibration data. At one compression rank, the compressed model runs about six
percent faster than the uncompressed model on a laptop GPU. At a slightly larger
rank, throughput matches baseline and perplexity goes up by about thirteen
percent. We spent a long time trying to figure out why it gets faster and the
answer turned out to be kernel fusion, not cache effects, though the cache story
is partially consistent with the data. The paper is honest about what is measured
and what is not.

\vspace{6pt}

\noindentPaper~II --- Geodesic Projection Pipeline.}
Extends the compression from three attention slots to all seven compressible
matrices in a transformer block, adds per-layer rank allocation, and describes
four adaptive mechanisms (MCR, GL(d) gauge, thermal control, online Oja updates).
Only the MCR mechanism has been tested, and it was a null result at the tested
rank. The remaining mechanisms are design specifications, not measurements. The
paper distinguishes clearly between what has been built and measured versus what
has been designed on paper.

\vspace{6pt}

\noindentPaper~III --- Speculative Decoding with Compression.}
Combines the compressed drafter from Paper~I with standard speculative decoding.
On SmolLM2-135M we measure a $1.53{\times}$ speedup at about thirty-nine percent
acceptance rate. The paper also documents a failure mode where greedy speculative
decoding breaks against instruct-tuned models and the two-line fix that resolves
it. One finding is negative: aggressive compression and speculative decoding do
not compose well, which contradicts the initial hypothesis.

\vspace{6pt}

\noindentPaper~IV --- Organic Training Theory.}
The theoretical backbone. Proposes modeling transformer latent space as a
Riemannian manifold with the Fisher information metric where inference becomes
approximate geodesic flow. Also introduces Geodesic Trajectory Caching (GTC),
which stores completed geodesics and serves new queries by a single linear
correction. Twelve of seventeen testable claims now have measured anchors across
three model scales. Five remain open.

\vspace{6pt}

\noindentPaper~V --- GRC Light Distillation.}
An optional distillation step that recovers some of the perplexity lost to
compression. On SmolLM2-135M, a small LoRA adapter ($r{=}8$) recovers over one
hundred percent of the PPL gap at $k{=}512$, meaning the distilled compressed
model beats the uncompressed baseline on perplexity. At aggressive compression
($k{=}256$) recovery is partial, and at conservative compression ($k{=}1024$) it
is near-complete. The protocol is designed so that disabling distillation
produces bit-identical output to Paper~I.

\vspace{6pt}

\noindentPaper~VI --- Task-Level Impact.}
Measures what actually happens to benchmark scores when you compress attention.
On SmolLM2-135M, knowledge recall tasks like MMLU hold up reasonably well down
to $k{=}512$, but reasoning tasks degrade sharply. An earlier measurement
claiming twenty-eight percent MMLU at $k{=}256$ was retracted after discovering
the binary lacked proper ChatML template support. The corrected measurement is
zero percent at that rank.

\vspace{6pt}

\noindentPaper~VII --- FFN Cluster Compression.}
Attempts structure-aware compression of the feed-forward network using column
clustering plus per-cluster SVD. Four-cluster compression reduces reconstruction
error by about twenty-three percent versus global SVD. The approach recovers some
of the information that global methods lose, but the overall gain is modest. The
paper includes a correction from an earlier overoptimistic prediction.

\vspace{6pt}

\noindentPaper~VIII --- GTC as RAG Alternative.}
Explores whether the geodesic trajectory cache from Paper~IV can serve as a
retrieval mechanism, comparing it conceptually against retrieval-augmented
generation. The computed speedup is large on paper (hundreds of times) but
cross-model generalization has not been measured. This is marked as pending
hardware availability.

\vspace{6pt}

\noindentPaper~IX --- The 106\% Anomaly as a General Phenomenon.}
Argues that the super-baseline speedup observed in Paper~I is not a quirk of
one GPU or one model. Derives the condition under which compressed kernels run
faster than uncompressed ones, predicts optimal compression ranks for ten GPU
types, and identifies six kernel classes beyond attention where the effect should
appear. Three kernel classes are confirmed by measurement; nine GPU predictions
remain open.

\vspace{6pt}

\noindentPaper~X --- Cross-Embedding Compatibility Index.}
Tests whether you can splice parts of different models together by aligning their
embedding spaces. Within the same model, splicing fails at low compression ranks
(zero compatible layers at $k{=}256$). Between different models trained on
similar data, splicing works at full rank when architecture matches. A working
chimera called MINSKAT (attention from one model, FFN from another) was built and
published. The CECI acronym was renamed mid-project from CMVB; some early drafts
may still use the old name.

\vspace{14pt}

\noindent Thank you again for reading. If you find errors, they are mine. If you
find something useful, the credit belongs to the open-weight model community that
made the weights available, the authors of the libraries I built on, and the
friends who read early drafts and told me which parts made no sense.

\vspace{8pt}

\begin{flushright}
William Ken Ohara Stewart

May 2026
\end{flushright}

\end{document}
